\documentclass[10pt,a4paper]{amsart}
\usepackage[margin=19mm]{geometry}
\usepackage{amsmath,amssymb,mathtools}
\usepackage[scr=rsfs,cal=euler]{mathalfa}
\usepackage{xfrac}
\usepackage[unicode,hidelinks,bookmarks=false]{hyperref}
\newcommand{\ie}{i.e.\ }
\newcommand{\cf}{cf.\ }
\newcommand{\resp}{respectively\ }
\newcommand{\Subsection}{\subsection}
\providecommand{\nobreakdash}{\nobreak\hbox{-}}
\setlength{\emergencystretch}{2em}
\multlinegap0pt
\usepackage{amssymb}
\usepackage{enumerate}
\usepackage{tikz}
\newtheorem{lem}{Lemma}
\def\GL{ \text{\rm GL} }
\title{On the source algebra equivalence class of blocks with cyclic defect groups, II}
\author{Gerhard Hiss, Caroline Lassueur}
\date{Source: arXiv:2502.09176v4 (CC BY 4.0)}
\begin{document}
\maketitle

\noindent\emph{Source and licence.} On the source algebra equivalence class of blocks with cyclic defect groups, II. Version arXiv:2502.09176v4 (CC BY 4.0), distributed by arXiv under the Creative Commons Attribution 4.0 International licence, http://creativecommons.org/licenses/by/4.0/, as recorded in the arXiv licence metadata for this exact version and shown on the abstract page https://arxiv.org/abs/2502.09176v4. Full licence notice: \texttt{licenses/arxiv-2502.09176-CC-BY-4.0.txt}.\par\smallskip

\noindent\textit{Editorial scope.} Counted unit: the article's lem MinimalPolynomialAndDeterminant (source0.tex lines 2080--2110). The article's complete original proof of it is delivered with it (source0.tex lines 2111--2155, one \texttt{proof} environment of 45 lines). No mathematics is written, repaired or re-derived: the statement and the proof are the article's own lines, in the article's own notation, and no retained line is substituted or re-flowed.  Retained uncounted as the source-local context the proof needs: lines 1426--1432.  Nothing was omitted from any retained span, so the package carries no omission record.  No retained line contains a citation, so the wrapper carries no bibliography and no bibliography entry.  No retained line contains a figure environment, an image-inclusion command or drawing code, so no image is delivered and the package declares no asset dependency.  Editorial formatting only, in addition: an AMS \texttt{amsart} preamble that restates the article's own declarations and macros used by the retained lines, namely the environments \texttt{lem} on separate counters, together with the standard packages those lines need.  The retained environments print in this document as Lemma 1 in order of appearance, whereas the article prints the same items with the numbering of its own full document.  The complete native source of the article as distributed in the arXiv e-print for this exact version is bundled unchanged as \texttt{sources/173227/source0.tex}.
\begin{lem}
\label{UnitaryDelta}
Suppose that $p \mid q + 1$. 
Let $\Delta$ be a monic irreducible polynomial over~$\mathbb{F}_{q^2}$ of odd 
degree~$e$, and let $\zeta \in \mathbb{F}$ be a root of~$\Delta$ of $p$-power 
order. Then $\Delta = \Delta^{\dagger}$.
\end{lem}

\begin{lem}
\label{MinimalPolynomialAndDeterminant}
Let $t \in G$ be semisimple and let~$\Gamma$ denote the minimal polynomial 
of~$t$ over $\mathbb{F}_{q^{\delta}}$. Let~$T$ denote a cyclic maximal torus 
of~$G$ of order $q^n - \varepsilon^n$.

{\rm (a)} If~$\Gamma$ is irreducible, the minimal polynomial of~$t^j$ is 
irreducible for all integers~$j$. (This does not use the hypothesis
that $p \mid q - \varepsilon$.)

{\rm (b)} Suppose that $t \in T$. Then~$\Gamma$ is irreducible, unless 
$\varepsilon = -1$, $n$~is even and $\Gamma = \Delta\Delta^{\dagger}$ for a 
monic irreducible polynomial $\Delta$ with $\Delta \neq \Delta^{\dagger}$.

{\rm (c)} Let~$\Delta$ be a monic irreducible polynomial 
over~$\mathbb{F}_{q^{\delta}}$, and let $\zeta \in \mathbb{F}$ be a root 
of~$\Delta$. If~$\zeta$ is of $p$-power order, then
$\deg(\Delta) = p^b$ with
$$
b = \begin{cases} 0, & \text{\ if\ } |\zeta| \leq p^a \\
                    a', & \text{\ if\ } |\zeta| = p^{a + a'} \text{\ with\ } a' \geq 0.
\end{cases}
$$

{\rm (d)} If $t \in T$ is of $p$-power order, then~$\Gamma$ is irreducible.

{\rm (e)} Suppose that~$\Gamma$ is irreducible and put $d := \deg( \Gamma )$.
Thus $n = n' d$ for some positive integer~$n'$. Let~$\zeta \in \mathbb{F}$ be 
an eigenvalue of~$t$. Then
$$\det(t) = \zeta^{n'(q^{\delta d} -1)/(q^{\delta}-1)}.$$
\end{lem}

\begin{proof}
(a) This is certainly well known. For convenience, we sketch a proof.
As~$\Gamma$ is irreducible, the subalgebra $\mathbb{F}_{q^{\delta}}[t]$ of the 
full matrix algebra is a field. It is, in particular, a finite extension 
of~$\mathbb{F}_{q^{\delta}}$. Thus~$t^j$ is algebraic 
over~$\mathbb{F}_{q^{\delta}}$, and so $\mathbb{F}_{q^{\delta}}[t^j]$ is a 
field. This implies the claim.

(b) Suppose that $\varepsilon = 1$ or that~$n$ is odd. Then~$T$ lies in a 
Coxeter torus of~$\GL_n( q^{\delta} )$. Hence~$t$ is a power of an element 
of~$\GL_n( q^{\delta} )$ acting irreducibly on~$V$. 
Thus~$\Gamma$ is irreducible by~(a). 

Now suppose that $\varepsilon = -1$ and that~$n = 2m$ is even. Then~$T$ is a
Coxeter torus of a Levi subgroup $L \cong \GL_m( q^2 )$ of~$G$. Now 
$V = V_1 \oplus V_2$ with $L$-invariant totally isotropic 
subspaces~$V_1$, $V_2$ of~$V$, such that~$L$ acts on~$V_1$ in 
the natural way, i.e.\ as $\GL_m( q^2 )$ acts on~$\mathbb{F}^m_{q^2}$, whereas 
the action of~$L$ on~$V_2$ is the natural action twisted by an automorphism 
of~$L$. Thus, by the first part of the proof, if~$\Gamma$ is reducible, we have 
$\Gamma = \Delta \Delta'$ with monic irreducible polynomials 
$\Delta \neq \Delta'$. As~$t$ lies in the unitary group, we must have 
$\Delta' = \Delta^{\dagger}$.

(c) Put $d = \deg( \Delta )$. If $|\zeta| \leq p^a$, then 
$\zeta \in \mathbb{F}_{q^{\delta}}$, and thus $d = 1$. 
Suppose then that $|\zeta|= p^{a + a'}$ for some non-negative integer~$a'$. 
Observe that the highest power of~$p$ dividing $q^{\delta} - 1$, respectively 
$q^{\delta p^{a'}} - 1$, equals~$p^a$, respectively $p^{a+a'}$. 
Hence
$d = [\mathbb{F}_{q^{\delta}}[\zeta]\colon\!\mathbb{F}_{q^{\delta}}] = p^{a'}$. 

(d) Suppose that~$\Gamma$ is reducible, and let $\zeta \in \mathbb{F}$ denote
a root of~$\Delta$, where~$\Delta$ is as in~(b). Then~$\zeta$ is of $p$-power 
order, and thus $\deg(\Delta)$ is odd by~(c). But then 
$\Delta = \Delta^{\dagger}$ by Lemma~\ref{UnitaryDelta}, a contradiction.

(e) As $\Gamma$ is irreducible,
the roots of~$\Gamma$ are
$$\zeta, \zeta^{q^{\delta}}, \zeta^{q^{\delta \cdot 2}}, \ldots ,
\zeta^{q^{\delta \cdot(d - 1)}},$$
and the product of these roots equals
$$\zeta^{(q^{\delta d} - 1)/(q^{\delta} - 1)}.$$
This proves our claim.
\end{proof}

\end{document}
